\scanpage{413}
\mspage{327}
% archive sheet number: 412

% 1–6
indiqué page 322. Ensuite, on corrige \unsure{ce|ces} choix, en remplaçant $\unsure{j|J}_\alpha$ par $\overline{j}_\alpha$, \deleted{\unsure{disons|alors}} \inserted{avec} $\rho:j_\alpha\to\overline{j}_\alpha$, et $u_\alpha$ par $\overline{u}_\alpha=\rho u_\alpha$, $\tau'_{\alpha,\alpha_0}$ par $\overline{\tau'}_{\alpha,\alpha_0}=\rho\tau'_{\alpha,\alpha_0}$, de façon à satisfaire aux contraintes \unsure{(}$3_\alpha$)$(\beta)$ de la page 318.

% 7–8
Cela achève la construction dans le cas où $\alpha=\alpha_0+1$. \deleted{\unsure{où}}

% 9
\uline{Cas 2 B} : \uline{$\alpha$ est un ordinal limite}.

% 10–13
Il faut donc construire les données 1) 2) 3) de la page 317 (à l'exclusion de la donnée 4), de façon à satisfaire à la seule contrainte $3_\alpha(\beta)$.

% 14–18
On commence à construire $i_\alpha$ et les $\tau_{\beta,\alpha}$ dans $I$. C'est possible, car $I_\alpha=\{\beta\mid<\beta\deleted{\unsure{\alpha}}<\alpha\}$ est de cardinal \deleted{$\unsure{\pi_\alpha}$} $\pi_\alpha<\unsure{\pi|\varpi}$ -- il suffit donc de supposer $I$ \unsure{grd|grand} devant les cardinaux \deleted{\unsure{petits}} $\pi'<\pi$,

% 19–20
On construit ensuite de \unsure{$=$} $j_\alpha$ \inserted{dans $J'$} \deleted{et} les $\tau'_{\beta,\alpha}$ $(\beta<\alpha)$, en utilisant \unsure{le|les}
